\section{From exact screens to the physical interface of M-theory}
\label{mdp4:sec:interfaz-teoria-m-rev8}
\index{M theory! physical interface}
\index{T-duality}
\index{string theory}
\index{CFT}

The first part constructed, in explicit mathematical domains, the reductions
\[
12\longrightarrow11\longrightarrow10,
\qquad
26\longrightarrow24,
\]
and the moment, winding, and radius involutions
\[
(m,w,R)\longleftrightarrow(w,m,R^{-1}).
\]
These results are not yet a physical theory of strings nor an action
of supergravity. They constitute the screens and equivalences that a
dimensional realization must respect.

The interface is typed via a map declared
\[
\mathcal R_M:\mathcal X_{\rm HMT}^{\rm enr}
\longrightarrow
\mathcal X_{11}^{\rm phys},
\]
whose codomain must carry, at a minimum, the metric field \(g_{MN}\), the three-form \(C_{MNP}\), the gravitino \(\psi_M\), its action, and the supersymmetry transformations. A specialization to strings further requires identification of the physical radius, tension, momentum and winding sectors, and the state space on which T-duality is realized. A conformal or CFT reading requires its own operator algebra, grading, and state--operator correspondence.

The positive connection with the exceptional branch is not a bare arrow
Monster\(\to\)M-theory. Both prolongations preserve data from the same
enriched state---phase, grading, orientation, sewing, and memory---
but employ distinct functors. In the exceptional branch, those data reach
\(V^\natural\), the Monster, and Moonshine; in the dimensional branch they feed
the screens and, only through \(\mathcal R_M\), a possible realization of
fields and strings.

\begin{definition}[Internal Admissibility of a Compactification]
\label{mdp4:def:admisibilidad-compactificacion-hmt-rev3}
A candidate compactification is represented by a triple
\[
 (x_\infty,\kappa,\mathcal U),
\]
where \(x_\infty\in\mathcal X_{\rm HMT}^{\rm enr}\) is a surviving
section,
\[
 \kappa:\mathcal X_{\rm HMT}^{\rm enr}\dashrightarrow\mathcal M
\]
is the possibly partial map into a declared moduli space
\(\mathcal M\), and
\(\mathcal U\) transports metric, orientation, register and boundary data.
We define
\[
 \operatorname{Adm}_{\rm HMT}(x_\infty,\kappa,\mathcal U)=1
\]
if the following four conditions hold:
\begin{enumerate}
 \item the section has a compatible extension to every depth;
 \item the declared dualities are compatible at the point considered:
       \[
        \kappa(\mathsf T_{\rm HMT}x_\infty)
        =\mathsf T_{\mathcal M}(\kappa(x_\infty)),
       \]
       when both sides are defined;
 \item the spectrum descends to the compactification quotient;
 \item the observables preserve their domains, codomains, and types under
       \(\mathcal U\).
\end{enumerate}
\end{definition}

This predicate is an internal admissibility criterion and not an automatic identification with the physical \emph{Swampland}. The distance, tower, weak-gravity, and absence-of-global-symmetries criteria, as well as the geometric conditions of a specific compactification, are subsequently evaluated on \(\operatorname{im}\kappa\) and require the corresponding metric maps, fields, and scales. The preceding definition fixes the interface that a realization must preserve; it does not by itself prove the existence of a physical compactification or a Swampland conjecture.

\begin{statusbox}[title={Scope of the interface}]
The screens, quotients, and T-dual involution are exact
mathematical results. The complete dictionary toward
\((g_{MN},C_{MNP},\psi_M)\), the action, supersymmetry, and a physical CFT is
a typed interface that must be verified in each realization. This boundary
does not diminish the mathematical result or authorize its presentation as an already
closed physical dynamics.
\end{statusbox}
